\documentclass{article}
\usepackage{amsmath}
\usepackage{mathtools}
\usepackage{amsfonts}
\usepackage{parskip}
\newcommand*\plus{+{}}
\newcommand*\boxSizeOfMax[1]{\makebox[\widthof{Max}][c]{#1}}
\setlength\parindent{0pt}

\title{A Career Development Plan Framework for Senior Professionals - DRAFT\\ [1em]
	\large Jonathan Erdmann \\ [1em]
}

\begin{document}

\maketitle

This document is the governing framework of a personal strategic career management system.
This framework defines the structure, purpose, and operating principles of the system as a means to address the lack of such
tools with my professional experience up to this point.
Every strategic decision, planning artifact, and review cadence developed towards my personal career management plan derives its 
authority from what is written here. When content and structure are in conflict, structure governs. When execution and strategy are in conflict,
this document provides the mechanism for resolution.

The system is modeled on general program management principles and adapted from a cursory overview of Hoshin Kanri strategy deployment.
It is designed to address the limitations in our existing development planning framework and is for one purpose:
sustained, intentional career development over a multi-year horizon, executed with the same rigor applied to 
consequential programs of record.

\section{Governing Principle}

Every element defined within the framework must trace to a parent.
An action item without an initiative, an initiative without an annual objective, or an annual objective without
a breakthrough goal either has no place in the plan or the plan has a gap.

\section{Structure}

The framework has three layers. Each layer has a distinct purpose, operates over a different time horizon,
and demands a different level of granularity. Information flows upward through the layers via the performance cadence.
Decisions flow downward such that strategy informs the operating model, and the operating model informs daily action.

\subsection{Layer 1: Strategic Foundation}

\textbf{Purpose}: To define where I am going and why I will thrive there.

This layer operates at the level of enduring truth and serves as an anchor for the strategy. It contains statements that 
should remain stable across the full planning horizon and require deliberate annual review before they are changed.
Language contained within this layer is the considered output of rigorous self-assessment and strategic analysis.
If elements cannot be defended, they should be eliminated.

\textbf{Leveling Guidance:} Layer 1 content is general enough to survive role changes, company changes, and market 
shifts without falsification. Its content is specific enough that a reader could not mistake it for belonging to a
different person. Generic statements that could describe any senior professional are insufficient.

\textbf{Contents:}

\begin{itemize}
\item \textbf{Personal Mission}: A single statement answering why \textit{this} direction, in my own words.
	The test is whether it remains true in five years regardless of which specific role or company I am in.
\item \textbf{Positioning Thesis}: The competitive argument for why specifically \textit{I} will thrive in my target functions, grounded in 
	specific evidence from my career history. Not what I want to be known for, but what I can actually demonstrate.
\item \textbf{Target Functions}: Defined roles with entry angles, competitive rationale, and sequencing logic.
	Specific enough to guide targeting decisions but not so narrow as to preclude adjacent opportunities.
\item \textbf{Known Gaps}: An honest, unvarnished assessment of what stands between current state and target state.
	Gaps are the inputs into Layer 2 with the intent to expose these areas to growth opportunities.
\item \textbf{Breakthrough Goals}: Three to five statements defining success over the 5-year horizon. Breakthrough Goals are written 
	as conditions, not activities. The test of a breakthrough goal is whether I would unambiguously know if I have achieved it.
\end{itemize}

\newpage

\subsection{Layer 2: Operating Model}

\textbf{Purpose}: Translate strategy into an executable, milestone-gated plan.

This layer is the living roadmap. This layer operates at the level of committed plans which are specific enough 
to be executed and flexible enough to be updated when Layer 3 surfaces a signal. Content here bridges the enduring truth 
of Layer 1 and the weekly reality of Layer 3. 

\textbf{Leveling Guidance:} Layer 2 content is specific enough that completion can be assessed without judgment.
Each initiative has a definition of done and each action item is discrete and ownable. Vague entries such as 
``build network'' are insufficient and should be decomposed until they are actionable.

\textbf{Element Hierarchy:}

\begin{itemize}
	\item \textbf{Breakthrough Goal} (Defined in Layer 1): The primary conditions which represent satisfaction of the strategy once met
	\item \textbf{Annual Objective}: What must be true by year end to be on track toward the breakthrough goal
	\item \textbf{Initiative}: The workstream or project that achieves the annual objective
	\item \textbf{Action Item}: The discrete, completable unit of work that advances the initiative
\end{itemize}

Each initiative has a definition of done and a target timeframe. Each phase of the plan has a named gate which
represents a condition (not just a date) that determines readiness to advance.

\subsection{Layer 3: Performance Cadence}

\textbf{Purpose}: Surface signals, drive accountability, and trigger updates to Layer 2.

This layer is the operating rhythm of the strategy and functions at the level of honest observation: what actually happened 
versus what was committed to, and what that pattern means. Content within this layer does not require polish as its a truthful 
log of commitments, completions, and reflections. Its value compounds over time as patterns become visible.

\textbf{Leveling Guidance:} Layer 3 entries are specific, dated, and without generalization. The weekly log records actual actions,
not intentions. The monthly reflection names specific misalignments, not general impressions. The quarterly review renders a 
judgment on whether the strategy is working as supported by observable evidence.

\newpage

\textbf{The Four Review Cycles:}

\begin{itemize}
	\item \textbf{Weekly}: Did I do what I committed to? Track actions committed versus actions completed without strategy discussion. 
		The purpose is building the habit of honoring commitments to myself and identifying where friction actually lives.
  \item \textbf{Monthly}: Are my activities aligned to my priorities? Look at the pattern across weeks. If a priority appears in 
	  	Layer 2 but no actions trace to it in the past month, that is a signal. Assess whether the misalignment is a motivation problem
		, a sequencing problem, or a strategy problem.
  \item \textbf{Quarterly}: Is my positioning working? This is where leading indicators are assessed. Are the right people finding me?
	  	Are conversations happening with the right functions? Are capability gaps closing? This is also where Layer 2 is formally 
		reviewed and updated if needed.
  \item \textbf{Annually}: Is the strategy itself still correct? Revisit Layer 1. The annual review asks whether the positioning thesis,
	  	target functions, and breakthrough goals still hold and revises the strategic foundation if they do not.
\end{itemize}

\section{The Correction Mechanism}

When the cadence surfaces a signal, a judgment is required before acting on it.

\textbf{Execution problem:} 
\begin{itemize}
	\item The strategy is correct but actions are not being completed, initiatives are stalled, or milestones 
are slipping. 
	\item Response: update Layer 2. Resequence, rescope, or recommit at the action item or initiative level.
\end{itemize}

\textbf{Strategy problem:}
\begin{itemize}
	\item Actions are being completed but leading indicators are not moving, or external conditions have materially changed.
	\item Response: flag for the annual review. Do not reactively update Layer 1 mid-cycle except under significant circumstances.
\end{itemize}

This distinction prevents two failure modes: the rigid plan that never adapts, and the reactive plan that is not consistent enough
to generate signal.

\newpage

\section{Populating Layer 1}

Establishing Layer 1 is the key activity in the strategy development effort and likely requires multiple sittings.
Each element requires a different type of thinking and a different input.

\textbf{Personal Mission}: Developed through directed reflection by examining what work genuinely energizes, what type of 
problem demands my best thinking, what the satisfaction of good work actually feels like, and why that matters personally. 
The mission statement is arrived at through honest interrogation of those questions until the answer is specific 
enough to be my own and stable enough to survive a long-term time horizon.

\textbf{Positioning Thesis:} Built from three components that must work together: differentiated capability, evidence base,
and target context. The thesis is drafted only after the competitive landscape of the target functions has been examined.
The thesis makes an additive argument and not comparative (e.g. I'm not defining myself by my perceived strengths/weaknesses of others).

\textbf{Target functions:} Defined with three elements: a functional description grounded in how the role actually operates 
, an entry angle that is honest about the credible path given my current positioning, and a time horizon that reflects 
sequencing logic rather than aspiration. Functions are pruned until the list is coherent and executable.

\textbf{Known Gaps:} Assessed in three stages corresponding to the three time horizons of the target functions.
Near-term gaps are assessed in detail across two dimensions: competency gaps - the analytical, financial, and professional 
capabilities required to perform in the target function, and positioning gaps - the visibility, credentialing, and network 
conditions required to be considered for it. Competency gaps are named first because execution of the mission depends on the 
quality of the work, not the quality of the positioning. Medium-term gaps are assessed directionally. Long-term 
gaps are noted at a high level and revisited at each annual review.

\textbf{Breakthrough Goals:} Written last, after the preceding elements have been established. They are written as conditions rather
than activities. Each must trace directly to the personal mission and positioning thesis, and each must have at least one annual 
objective in Layer 2 that traces back to it.

\newpage

\section{Populating Layer 2}

\textbf{Annual Objectives}: Derived by working backward from each breakthrough goal. The question is: what must be true at the end of this 
year for the breakthrough goal to remain achievable on its current trajectory? Annual objectives are written as conditions rather than 
activities, and are specific enough that progress can be assessed at year end without judgment. Each breakthrough goal may have one 
or more annual objectives, and their target years reflect the sequencing logic of the overall plan.

\textbf{Initiatives:} Derived by identifying the specific workstreams that achieve each annual objective. Each initiative has a clear 
definition of done and a target timeframe. Initiatives are the primary working unit of Layer 2. They are specific enough to be executed 
but scoped broadly enough that they represent a coherent body of work rather than a single task.

\textbf{Action Items:} are derived by decomposing each initiative into discrete, completable units of work. The test for an action item is that 
its completion can be assessed without ambiguity. Action items are the lowest node in the hierarchy and the primary interface of the
Layer 3 weekly cadence. They are added and updated continuously as the plan is executed, not defined exhaustively upfront.

The phase-gate structure organizes initiatives into phases, each with a named gate condition. A gate condition itself is not a date itself,
but rather a state of the world that must be true before the next phase begins. Gate conditions are derived from the breakthrough goals and 
objectives that anchor each phase. Advancement through phases is based on satisfaction of initiatives and not scheduled-based. Initiative 
detail for future phases is defined progressively at each annual review as the operating context becomes clearer and prior gate conditions are 
cleared or in sight.

%\newpage
%
%\section{Relationship to Existing Frameworks}
%
%This framework draws deliberately from three sources.
%
%\textbf{Hoshin Kanri} provides the core logic: A strategic foundation that stays stable, an operating model that executes it, and a cadence that surfaces feedback and drives course correction.
%The catchball concept which features upward flowing information and downward flowing decisions is the structural backbone.
%
%\textbf{Program Management} provides the phase-gate structure and the discipline of defining done before beginning.
%
%\textbf{The Work Breakdown Structure} provides the element hierarchy in Layer 2. Every action item traces to an initiative, every initiative to an annual objective, every annual objective to a breakthrough goal.

%\section{What This Framework Is Not}
%
%It is not a to-do list. Action items are the lowest node in a strategic hierarchy, not the starting point.
%
%It is not a goal-setting exercise. Breakthrough goals without an operating model and a cadence are wishes.
%
%It is not static. The correction mechanism is not a failure mode --- it is the intended behavior of a living system.
%
%It is not a performance evaluation. The cadence is a tool for self-directed accountability, not judgment.

\end{document}
